\begin{proof}[Proof of \cref{obj:lem:positive-ratio}]
\begin{enumerate}
\item Fix \(x\in\mathcal X\). Define the centered latent dose and stratum probability by
\[
T:=A-a_0,\qquad p_x:=P(X=x).
\]
Choose a measurable conditional density \(g_x:\mathbb R\to\mathbb R\) supplied by \cref{obj:ass:weak-design}, so that
\[
\mathcal L_P(T\mid X=x)(dt)
=
g_x(t)\mathbf1_{[-1/2,1/2]}(t)\,dt,
\qquad
c_{\mathrm{lo}}|t|^\kappa\le g_x(t)\le c_{\mathrm{hi}}|t|^\kappa
\]
for Lebesgue-almost every \(t\in[-1/2,1/2]\). To express the mean on the whole real line, define
\[
\bar\mu_x:\mathbb R\to\mathbb R,
\qquad
\bar\mu_x(a):=\mu_x\!\left(\min\{1,\max\{0,a\}\}\right).
\]
This extension is measurable by the Hölder condition, and \cref{obj:ass:mean-range} gives
\[
0\le\bar\mu_x(a)\le1\qquad(a\in\mathbb R).
\]
Define the latent weighted integrals
\[
L_{x,q}:=\int_{-1/2}^{1/2}g_x(t)q(t)\,dt,
\qquad
N_{x,q}:=\int_{-1/2}^{1/2}
g_x(t)\bar\mu_x(a_0+t)q(t)\,dt.
\]
Multiplying the design envelopes by the nonnegative weight gives
\[
0\le c_{\mathrm{lo}}|t|^\kappa q(t)
\le g_x(t)q(t)
\le c_{\mathrm{hi}}|t|^\kappa q(t)
\]
almost everywhere on the integration interval. The upper bound and the finite base moment ensure integrability of \(g_xq\); boundedness of \(\bar\mu_x\) then ensures integrability of the numerator integrand. Integrating the lower bound yields
\[
L_{x,q}\ge c_{\mathrm{lo}} I_\kappa(q)>0.
\]
Also, \cref{obj:ass:stratum-overlap} supplies
\[
p_x\ge p_{\min}>0.
\]
% lean: CausalSmith.Stat.NoisydoseWeakdesignTransition.latent_weight_lower

\item Define the normalized stratum law and the observed denominator by
\[
P_x:=\frac{1}{p_x}P|_{\{X=x\}},
\qquad
D_{x,q}:=E_P[\mathbf1\{X=x\}\ell_q(V)].
\]
Here \(P_x\) is a probability law on the structural coordinate space. Restriction and normalization preserve the assumed absolute integrability of \(\ell_q(V)\).

The error independence in \cref{obj:ass:error-independence} gives, for Borel sets \(C_Z,C_T\subseteq\mathbb R\),
\[
\begin{aligned}
P_x(Z\in C_Z,T\in C_T)
&=\frac{P(Z\in C_Z,X=x,T\in C_T)}{p_x}\\
&=P(Z\in C_Z)\frac{P(X=x,T\in C_T)}{p_x}\\
&=P(Z\in C_Z)P_x(T\in C_T).
\end{aligned}
\]
Together with \cref{obj:ass:gaussian-channel}, this identifies the joint law as
\[
\mathcal L_{P_x}(Z,T)
=
N(0,1)\otimes\mathcal L_{P_x}(T).
\]
% lean: CausalSmith.Stat.NoisydoseWeakdesignTransition.error_centered_product_stratum

Since \(V=T+\sigma Z\), Fubini's theorem, justified by absolute integrability, and the pointwise Gaussian-inverse identity give
\[
\begin{aligned}
\frac{D_{x,q}}{p_x}
&=E_{P_x}[\ell_q(T+\sigma Z)]\\
&=\int_{-1/2}^{1/2}
g_x(t)E_Z[\ell_q(t+\sigma Z)]\,dt\\
&=\int_{-1/2}^{1/2}g_x(t)q(t)\,dt
=L_{x,q}.
\end{aligned}
\]
Consequently,
\[
D_{x,q}=p_xL_{x,q}.
\]
% lean: CausalSmith.Stat.NoisydoseWeakdesignTransition.populationD_eq_latent

\item Define the observed numerator by
\[
M_{x,q}:=E_P[\mathbf1\{X=x\}Y\ell_q(V)].
\]
By \cref{obj:lem:realized-dose-mean},
\[
0\le Y\le1\quad P\text{-almost surely},
\qquad
E_P[Y\mid A,X,Z]=\mu_X(A)\quad P\text{-almost surely}.
\]
Thus \(M_{x,q}\) is finite. The multiplier
\(\mathbf1\{X=x\}\ell_q(V)\) is measurable with respect to \((A,X,Z)\) and integrable. Pulling it through conditional expectation gives
\[
\begin{aligned}
M_{x,q}
&=E_P\!\left[
E_P[\mathbf1\{X=x\}Y\ell_q(V)\mid A,X,Z]
\right]\\
&=E_P\!\left[
\mathbf1\{X=x\}\ell_q(V)E_P[Y\mid A,X,Z]
\right]\\
&=E_P[\mathbf1\{X=x\}\mu_x(A)\ell_q(V)]\\
&=E_P[\mathbf1\{X=x\}\bar\mu_x(a_0+T)\ell_q(T+\sigma Z)].
\end{aligned}
\]
The last equality uses the latent support in \cref{obj:ass:latent-support}, on which \(a_0+T=A\in[0,1]\).
% lean: CausalSmith.Stat.NoisydoseWeakdesignTransition.integrable_test_mean

The product \(\bar\mu_x(a_0+T)\ell_q(T+\sigma Z)\) is absolutely integrable because \(0\le\bar\mu_x\le1\). Applying the product-law integration from the preceding step with this bounded multiplier yields
\[
\begin{aligned}
\frac{M_{x,q}}{p_x}
&=E_{P_x}[\bar\mu_x(a_0+T)\ell_q(T+\sigma Z)]\\
&=\int_{-1/2}^{1/2}
g_x(t)\bar\mu_x(a_0+t)
E_Z[\ell_q(t+\sigma Z)]\,dt\\
&=\int_{-1/2}^{1/2}
g_x(t)\bar\mu_x(a_0+t)q(t)\,dt
=N_{x,q}.
\end{aligned}
\]
Hence
\[
M_{x,q}=p_xN_{x,q}.
\]
% lean: CausalSmith.Stat.NoisydoseWeakdesignTransition.populationM_eq_latent

\item The positivity established above implies
\[
D_{x,q}=p_xL_{x,q}>0.
\]
Moreover, multiplying the lower bounds for \(p_x\) and \(L_{x,q}\) gives
\[
D_{x,q}
\ge p_{\min}c_{\mathrm{lo}}I_\kappa(q)
=b_q(K).
\]
% lean: CausalSmith.Stat.NoisydoseWeakdesignTransition.positive_ratio

\item We next bound the centered latent numerator. The parameter assumptions give \(\beta>0\) and \(\kappa\ge0\). On \([-1/2,1/2]\),
\[
0\le |t|^\beta\le1,
\qquad
|t|^{\kappa+\beta}q(t)
=
\bigl(|t|^\kappa q(t)\bigr)|t|^\beta.
\]
Thus the higher-moment integrand is integrable, as the product of the integrable base-moment integrand and a bounded measurable function.
% lean: CausalSmith.Stat.NoisydoseWeakdesignTransition.higher_weightMoment_integrable

The integrability already established permits subtraction inside the integral:
\[
N_{x,q}-\mu_x(a_0)L_{x,q}
=
\int_{-1/2}^{1/2}
g_x(t)q(t)\bigl(\bar\mu_x(a_0+t)-\mu_x(a_0)\bigr)\,dt.
\]
For every \(t\in[-1/2,1/2]\), both doses belong to \([0,1]\), and \cref{obj:ass:holder-mean} gives
\[
a_0+t\in[0,1],
\qquad
\bar\mu_x(a_0+t)=\mu_x(a_0+t),
\qquad
|\bar\mu_x(a_0+t)-\mu_x(a_0)|\le |t|^\beta.
\]
Using \(g_x(t)q(t)\ge0\) and the upper design envelope, we obtain almost everywhere
\[
\begin{aligned}
\left|g_x(t)q(t)
\bigl(\bar\mu_x(a_0+t)-\mu_x(a_0)\bigr)\right|
&=g_x(t)q(t)
|\bar\mu_x(a_0+t)-\mu_x(a_0)|\\
&\le c_{\mathrm{hi}}|t|^\kappa q(t)|t|^\beta\\
&=c_{\mathrm{hi}}|t|^{\kappa+\beta}q(t).
\end{aligned}
\]
The absolute value of an integral is bounded by the integral of the absolute value, so
\[
\begin{aligned}
|N_{x,q}-\mu_x(a_0)L_{x,q}|
&\le
\int_{-1/2}^{1/2}
\left|g_x(t)q(t)
\bigl(\bar\mu_x(a_0+t)-\mu_x(a_0)\bigr)\right|\,dt\\
&\le c_{\mathrm{hi}}I_{\kappa+\beta}(q).
\end{aligned}
\]
% lean: CausalSmith.Stat.NoisydoseWeakdesignTransition.latent_weight_centered_bias

\item Nonnegativity of the higher-moment integrand gives
\[
I_{\kappa+\beta}(q)\ge0.
\]
Since \(p_x>0\) and \(L_{x,q}>0\), the observed numerator and denominator identities imply
\[
m_{x,q}
=\frac{p_xN_{x,q}}{p_xL_{x,q}}
=\frac{N_{x,q}}{L_{x,q}},
\qquad
m_{x,q}-\mu_x(a_0)
=
\frac{N_{x,q}-\mu_x(a_0)L_{x,q}}{L_{x,q}}.
\]
Therefore,
\[
\begin{aligned}
|m_{x,q}-\mu_x(a_0)|
&=\frac{|N_{x,q}-\mu_x(a_0)L_{x,q}|}{L_{x,q}}\\
&\le\frac{c_{\mathrm{hi}}I_{\kappa+\beta}(q)}{L_{x,q}}\\
&\le\frac{c_{\mathrm{hi}}I_{\kappa+\beta}(q)}{c_{\mathrm{lo}}I_\kappa(q)}\\
&=\frac{c_{\mathrm{hi}}}{c_{\mathrm{lo}}}
\frac{I_{\kappa+\beta}(q)}{I_\kappa(q)}
=B_q(K).
\end{aligned}
\]
The argument applies to every \(x\in\mathcal X\). Gaussian inversion was used directly throughout the stated range of \(\sigma\), including \(\sigma=0\), so the same conclusions hold at that endpoint.
% lean: CausalSmith.Stat.NoisydoseWeakdesignTransition.positive_ratio
\end{enumerate}
\end{proof}
